\section{Limitaciones y matriz de riesgos técnicos}

Las principales limitaciones son: captura histórica sin fecha, objetivo dinámico e imperfecto,
muestra balanceada artificialmente, variables contextuales ausentes, análisis observacional sin
causalidad y falta de una evaluación fuera de muestra. También se desconoce si el conjunto cubre de
forma homogénea regiones, periodos, idiomas y niveles de exposición.

\begin{table}[H]
\centering
\caption{Matriz de riesgos ajustada a la evidencia disponible.}
\label{tab:riesgos}
\small
\begin{tabularx}{\textwidth}{p{3.1cm}p{1.7cm}p{2cm}X}
\toprule
Riesgo & Impacto & Probabilidad & Mitigación \\
\midrule
Leakage o memorización por \code{track_id} & Alto & Alta si se incluye & Excluir del diseño; conservar para trazabilidad \\
Memorización por artista & Alto & Alta & Comparar partición aleatoria con grupos por artista; controlar encoding \\
Deriva temporal & Alto & Alta & Obtener fechas, validar temporalmente y monitorear desempeño \\
Muestra no representativa & Alto & Alta & Limitar generalización; documentar cobertura y reponderar solo con marco válido \\
Variables omitidas & Alto & Alta & Incorporar contexto trazable; no interpretar causalmente \\
Multicolinealidad de audio & Medio & Alta & Escalar, regularizar y diagnosticar coeficientes \\
Géneros multietiqueta & Medio & Alta & Multi-hot, soporte explícito y métricas no excluyentes \\
Extremos válidos confundidos con errores & Medio & Media & Validación puntual y análisis de sensibilidad \\
Desempeño desigual por segmentos & Alto & Media & Evaluar errores por género, explicit y artistas vistos/no vistos \\
Sobreconfianza de negocio & Alto & Media & Revisión humana, documentación y límites de uso \\
\bottomrule
\end{tabularx}
\end{table}

Probabilidad e impacto son evaluaciones cualitativas para priorizar controles, no frecuencias
estimadas. Deben revisarse cuando existan particiones, modelo, datos temporales y contexto de uso
concreto.
